% Options for packages loaded elsewhere
\PassOptionsToPackage{unicode}{hyperref}
\PassOptionsToPackage{hyphens}{url}
\documentclass[
]{article}
\usepackage{xcolor}
\usepackage[margin=1in]{geometry}
\usepackage{amsmath,amssymb}
\setcounter{secnumdepth}{-\maxdimen} % remove section numbering
\usepackage{iftex}
\ifPDFTeX
  \usepackage[T1]{fontenc}
  \usepackage[utf8]{inputenc}
  \usepackage{textcomp} % provide euro and other symbols
\else % if luatex or xetex
  \usepackage{unicode-math} % this also loads fontspec
  \defaultfontfeatures{Scale=MatchLowercase}
  \defaultfontfeatures[\rmfamily]{Ligatures=TeX,Scale=1}
\fi
\usepackage{lmodern}
\ifPDFTeX\else
  % xetex/luatex font selection
\fi
% Use upquote if available, for straight quotes in verbatim environments
\IfFileExists{upquote.sty}{\usepackage{upquote}}{}
\IfFileExists{microtype.sty}{% use microtype if available
  \usepackage[]{microtype}
  \UseMicrotypeSet[protrusion]{basicmath} % disable protrusion for tt fonts
}{}
\makeatletter
\@ifundefined{KOMAClassName}{% if non-KOMA class
  \IfFileExists{parskip.sty}{%
    \usepackage{parskip}
  }{% else
    \setlength{\parindent}{0pt}
    \setlength{\parskip}{6pt plus 2pt minus 1pt}}
}{% if KOMA class
  \KOMAoptions{parskip=half}}
\makeatother
\usepackage{color}
\usepackage{fancyvrb}
\newcommand{\VerbBar}{|}
\newcommand{\VERB}{\Verb[commandchars=\\\{\}]}
\DefineVerbatimEnvironment{Highlighting}{Verbatim}{commandchars=\\\{\}}
% Add ',fontsize=\small' for more characters per line
\usepackage{framed}
\definecolor{shadecolor}{RGB}{248,248,248}
\newenvironment{Shaded}{\begin{snugshade}}{\end{snugshade}}
\newcommand{\AlertTok}[1]{\textcolor[rgb]{0.94,0.16,0.16}{#1}}
\newcommand{\AnnotationTok}[1]{\textcolor[rgb]{0.56,0.35,0.01}{\textbf{\textit{#1}}}}
\newcommand{\AttributeTok}[1]{\textcolor[rgb]{0.13,0.29,0.53}{#1}}
\newcommand{\BaseNTok}[1]{\textcolor[rgb]{0.00,0.00,0.81}{#1}}
\newcommand{\BuiltInTok}[1]{#1}
\newcommand{\CharTok}[1]{\textcolor[rgb]{0.31,0.60,0.02}{#1}}
\newcommand{\CommentTok}[1]{\textcolor[rgb]{0.56,0.35,0.01}{\textit{#1}}}
\newcommand{\CommentVarTok}[1]{\textcolor[rgb]{0.56,0.35,0.01}{\textbf{\textit{#1}}}}
\newcommand{\ConstantTok}[1]{\textcolor[rgb]{0.56,0.35,0.01}{#1}}
\newcommand{\ControlFlowTok}[1]{\textcolor[rgb]{0.13,0.29,0.53}{\textbf{#1}}}
\newcommand{\DataTypeTok}[1]{\textcolor[rgb]{0.13,0.29,0.53}{#1}}
\newcommand{\DecValTok}[1]{\textcolor[rgb]{0.00,0.00,0.81}{#1}}
\newcommand{\DocumentationTok}[1]{\textcolor[rgb]{0.56,0.35,0.01}{\textbf{\textit{#1}}}}
\newcommand{\ErrorTok}[1]{\textcolor[rgb]{0.64,0.00,0.00}{\textbf{#1}}}
\newcommand{\ExtensionTok}[1]{#1}
\newcommand{\FloatTok}[1]{\textcolor[rgb]{0.00,0.00,0.81}{#1}}
\newcommand{\FunctionTok}[1]{\textcolor[rgb]{0.13,0.29,0.53}{\textbf{#1}}}
\newcommand{\ImportTok}[1]{#1}
\newcommand{\InformationTok}[1]{\textcolor[rgb]{0.56,0.35,0.01}{\textbf{\textit{#1}}}}
\newcommand{\KeywordTok}[1]{\textcolor[rgb]{0.13,0.29,0.53}{\textbf{#1}}}
\newcommand{\NormalTok}[1]{#1}
\newcommand{\OperatorTok}[1]{\textcolor[rgb]{0.81,0.36,0.00}{\textbf{#1}}}
\newcommand{\OtherTok}[1]{\textcolor[rgb]{0.56,0.35,0.01}{#1}}
\newcommand{\PreprocessorTok}[1]{\textcolor[rgb]{0.56,0.35,0.01}{\textit{#1}}}
\newcommand{\RegionMarkerTok}[1]{#1}
\newcommand{\SpecialCharTok}[1]{\textcolor[rgb]{0.81,0.36,0.00}{\textbf{#1}}}
\newcommand{\SpecialStringTok}[1]{\textcolor[rgb]{0.31,0.60,0.02}{#1}}
\newcommand{\StringTok}[1]{\textcolor[rgb]{0.31,0.60,0.02}{#1}}
\newcommand{\VariableTok}[1]{\textcolor[rgb]{0.00,0.00,0.00}{#1}}
\newcommand{\VerbatimStringTok}[1]{\textcolor[rgb]{0.31,0.60,0.02}{#1}}
\newcommand{\WarningTok}[1]{\textcolor[rgb]{0.56,0.35,0.01}{\textbf{\textit{#1}}}}
\usepackage{graphicx}
\makeatletter
\newsavebox\pandoc@box
\newcommand*\pandocbounded[1]{% scales image to fit in text height/width
  \sbox\pandoc@box{#1}%
  \Gscale@div\@tempa{\textheight}{\dimexpr\ht\pandoc@box+\dp\pandoc@box\relax}%
  \Gscale@div\@tempb{\linewidth}{\wd\pandoc@box}%
  \ifdim\@tempb\p@<\@tempa\p@\let\@tempa\@tempb\fi% select the smaller of both
  \ifdim\@tempa\p@<\p@\scalebox{\@tempa}{\usebox\pandoc@box}%
  \else\usebox{\pandoc@box}%
  \fi%
}
% Set default figure placement to htbp
\def\fps@figure{htbp}
\makeatother
\setlength{\emergencystretch}{3em} % prevent overfull lines
\providecommand{\tightlist}{%
  \setlength{\itemsep}{0pt}\setlength{\parskip}{0pt}}
\usepackage{bookmark}
\IfFileExists{xurl.sty}{\usepackage{xurl}}{} % add URL line breaks if available
\urlstyle{same}
\hypersetup{
  pdftitle={The change of affordability and availability in the housing market for young adults in the Netherlands from 2005 until 2025},
  pdfauthor={Björn Hildering 2887258, Tim Knegt 2882977, Willem van Beuningen 2787054, Robin Versnel 2862644, Seeger Tromp, Jaydane Levant 2851422},
  hidelinks,
  pdfcreator={LaTeX via pandoc}}

\title{The change of affordability and availability in the housing
market for young adults in the Netherlands from 2005 until 2025}
\author{Björn Hildering 2887258, Tim Knegt 2882977, Willem van Beuningen
2787054, Robin Versnel 2862644, Seeger Tromp, Jaydane Levant 2851422}
\date{2026-06-18}

\begin{document}
\maketitle

\section{Set-up your environment}\label{set-up-your-environment}

\begin{Shaded}
\begin{Highlighting}[]
\CommentTok{\# install.packages("tidyverse")}
\CommentTok{\# install.packages("yaml")}
\CommentTok{\# install.packages("rmarkdown")}
\CommentTok{\# install.packages("readr")}
\CommentTok{\# install.packages("dplyr")}
\CommentTok{\# install.packages("ggplot2")}
\CommentTok{\# install.packages(c(}
\CommentTok{\#   "geodata", "terra", }
\CommentTok{\#   "janitor", "scales", "patchwork"))}
\CommentTok{\# install.packages("tinytex")}
\end{Highlighting}
\end{Shaded}

\begin{Shaded}
\begin{Highlighting}[]
\FunctionTok{library}\NormalTok{(tidyverse)}
\end{Highlighting}
\end{Shaded}

\begin{verbatim}
## -- Attaching core tidyverse packages ------------------------ tidyverse 2.0.0 --
## v dplyr     1.2.1     v readr     2.2.0
## v forcats   1.0.1     v stringr   1.6.0
## v ggplot2   4.0.3     v tibble    3.3.1
## v lubridate 1.9.5     v tidyr     1.3.2
## v purrr     1.2.2     
## -- Conflicts ------------------------------------------ tidyverse_conflicts() --
## x dplyr::filter() masks stats::filter()
## x dplyr::lag()    masks stats::lag()
## i Use the conflicted package (<http://conflicted.r-lib.org/>) to force all conflicts to become errors
\end{verbatim}

\begin{Shaded}
\begin{Highlighting}[]
\FunctionTok{library}\NormalTok{(rmarkdown)}
\FunctionTok{library}\NormalTok{(yaml)}
\FunctionTok{library}\NormalTok{(readr)}
\FunctionTok{library}\NormalTok{(dplyr)}
\FunctionTok{library}\NormalTok{(ggplot2)}
\FunctionTok{library}\NormalTok{(geodata)}
\end{Highlighting}
\end{Shaded}

\begin{verbatim}
## Loading required package: terra
## terra 1.9.27
## 
## Attaching package: 'terra'
## 
## The following object is masked from 'package:tidyr':
## 
##     extract
\end{verbatim}

\begin{Shaded}
\begin{Highlighting}[]
\FunctionTok{library}\NormalTok{(terra)}
\FunctionTok{library}\NormalTok{(janitor)}
\end{Highlighting}
\end{Shaded}

\begin{verbatim}
## 
## Attaching package: 'janitor'
## 
## The following object is masked from 'package:terra':
## 
##     crosstab
## 
## The following objects are masked from 'package:stats':
## 
##     chisq.test, fisher.test
\end{verbatim}

\begin{Shaded}
\begin{Highlighting}[]
\FunctionTok{library}\NormalTok{(scales)}
\end{Highlighting}
\end{Shaded}

\begin{verbatim}
## 
## Attaching package: 'scales'
## 
## The following object is masked from 'package:terra':
## 
##     rescale
## 
## The following object is masked from 'package:purrr':
## 
##     discard
## 
## The following object is masked from 'package:readr':
## 
##     col_factor
\end{verbatim}

\begin{Shaded}
\begin{Highlighting}[]
\FunctionTok{library}\NormalTok{(patchwork)}
\end{Highlighting}
\end{Shaded}

\begin{verbatim}
## 
## Attaching package: 'patchwork'
## 
## The following object is masked from 'package:terra':
## 
##     area
\end{verbatim}

\begin{Shaded}
\begin{Highlighting}[]
\FunctionTok{library}\NormalTok{(tinytex)}
\end{Highlighting}
\end{Shaded}

Tutorial group 1 C. Schouwenaar \# Chapter 1: Identify a Social Problem
The Dutch housing market has changed significantly over the past two
decades. For young adults, finding suitable housing has become
increasingly difficult due to rising house prices, higher rents, and a
continuing housing shortage (Nederlands Jeugdinstituut, n.d.). As a
result, many young people struggle to enter the housing market.

This research examines how housing affordability and availability for
young adults in the Netherlands changed between 2005 and 2025. It
focuses on trends in house prices, rental costs, incomes and housing
shortages to understand the main developments during this period.

\subsection{1.1 Describe the Social
Problem}\label{describe-the-social-problem}

The Dutch housing market has become increasingly inaccessible for young
adults, making the move toward independent living far more difficult.
This is socially relevant because stable and affordable housing is a
foundation for financial independence, mental well‑being, and equal life
opportunities. When young people cannot secure a home, they delay major
life steps such as studying, working full‑time, or forming
households---effects that ripple through society. The OECD notes that
the Netherlands is among the countries where young adults face the
greatest barriers to affordable housing due to rising prices and
persistent shortages (OECD, 2023).

A combination of sharply increasing house prices, higher rents, and
limited supply has created a structural imbalance between demand and
availability. These pressures fall disproportionately on first‑time
buyers and young renters, widening generational inequality: older
households benefit from accumulated housing wealth, while younger
cohorts struggle to enter the market at all.

\section{Chapter 2: Data Sourcing}\label{chapter-2-data-sourcing}

For our analysis we use 5 datasets:

\subsection{2.1 Load in the data}\label{load-in-the-data}

\begin{Shaded}
\begin{Highlighting}[]
\CommentTok{\#Here we load in the datasets. The last one is the only one we immediately saved to data 2.0 since it was already good for use}
\NormalTok{inflatie }\OtherTok{\textless{}{-}} \FunctionTok{read.csv}\NormalTok{(}\StringTok{"data/Inflatie.csv"}\NormalTok{)}
\NormalTok{welvaart }\OtherTok{\textless{}{-}} \FunctionTok{read.csv}\NormalTok{(}\StringTok{"data/Welvaart\_van\_huishoudens\_\_kerncijfers.csv"}\NormalTok{)}
\NormalTok{bevolkingsgroei }\OtherTok{\textless{}{-}}\FunctionTok{read.csv}\NormalTok{(}\StringTok{"data/Bevolkingsgroei\_opbouw\_2005\_2019.csv"}\NormalTok{)}
\NormalTok{huizenvoorraad }\OtherTok{\textless{}{-}}\FunctionTok{read.csv}\NormalTok{(}\StringTok{"data/Huizenprijzen\_Huizenvoorraad.csv"}\NormalTok{)}
\NormalTok{Bestaande\_koopwoningen\_verkoopprijzen\_woningtype }\OtherTok{\textless{}{-}} \FunctionTok{read.csv}\NormalTok{(}\StringTok{"data/Bestaande\_koopwoningen\_verkoopprijzen\_woningtype.csv"}\NormalTok{)}
\NormalTok{brutoinkomen }\OtherTok{\textless{}{-}} \FunctionTok{read\_csv}\NormalTok{(}\StringTok{"data 2.0/bruto\_en\_besteedbaar\_inkomen\_leeftijdsgroep\_nieuw.csv"}\NormalTok{)}
\end{Highlighting}
\end{Shaded}

\begin{verbatim}
## Rows: 14 Columns: 9
## -- Column specification --------------------------------------------------------
## Delimiter: ","
## chr (1): Kenmerken_van_huishoudens
## dbl (8): Jaartal, Particuliere_huishoudens_x1000, Totaal_bruto_inkomen_mln.e...
## 
## i Use `spec()` to retrieve the full column specification for this data.
## i Specify the column types or set `show_col_types = FALSE` to quiet this message.
\end{verbatim}

\subsection{2.2 Short summary of the
datasets}\label{short-summary-of-the-datasets}

We chose these datasets since they give us valuable information about
affordability and availability of housing in the Netherlands.

\subsubsection{2.3 Description and Data
Collection}\label{description-and-data-collection}

The variables included in our project are \textbf{quantitative} in
nature. They consist of continuous economic indicators---such as
monetary values (income, house prices) and percentages (inflation)---as
well as discrete metrics tracking physical counts (population size,
housing stock). Categorical (nominal) variables such as
\texttt{Provincies} and \texttt{Woningtype} are utilized to allow for
spatial and structural segmentation.

Unlike survey data which relies on self-reported interviews, these
datasets are sourced entirely from \textbf{administrative registries}
compiled by Statistics Netherlands (CBS). Specifically, the
socio-economic and income data is derived from the Dutch Tax Authority
(Belastingdienst), while the housing metrics are registered through the
Municipal Base Registry for Housing (BAG) and the Dutch Land Registry
(Kadaster). This administrative nature ensures high data coverage and
eliminates traditional interview biases such as recall or non-response
bias.

\subsubsection{2.4 Stated limitations of this
data}\label{stated-limitations-of-this-data}

While these administrative datasets are highly reliable, we face several
limitations when merging and analyzing them to study young adults:

\textbf{Inconsistent Demographic Definitions:} The datasets categorize
``young adults'' differently. The income datasets slice data by the age
of the primary wage earner (\texttt{leeftijd\_kostwinner}), whereas the
demographic dataset tracks individuals aged 20 to 45
(\texttt{jong\_volwassenen\_20\_tot\_45\_jaar}). This structural
mismatch limits our ability to perfectly isolate a single, uniform age
bracket. \textbf{Temporal Truncation:} The demographic tracking for
population build-up (\texttt{Bevolkingsgroei\_opbouw\_2005\_2019.csv})
is currently limited to the 2005--2019 period within our data folder.
This creates a data gap for the final years of our 2005--2025 analysis
window. \textbf{Level of Aggregation:} Because the data is highly
aggregated by the CBS to protect privacy, we can only join these files
on the macro-level using the temporal dimension
(\texttt{Jaartal}/\texttt{Jaren}) and the regional dimension
(\texttt{Provincies}). We lack micro-data, meaning we cannot track
individual household trajectories or specific cross-sections over time.

\section{Chapter 3 - Quantifying}\label{chapter-3---quantifying}

\subsection{3.1 Data cleaning}\label{data-cleaning}

In this section we will clean up our datasets. We will show the progress
only ones, because it its quite repetitive.

Most of the collums have vague or ambiguous names, therefore we found it
necessary to alter these using the rename() function to make the table
more clear and easier to work with. To add to that, not all collums are
useful so we leave out some information to shrink the tables and only
keep the useful information. For this we use the function select(). When
the dataset is ready for use we will save it under a new name.
\texttt{\{r.content-hidden\ when-format="pdf"\}\ \#Here\ we\ modify\ the\ dataset\ Inflatie\ inflatie\ \textless{}-\ inflatie\ \%\textgreater{}\%\ \ \ rename(Jaartal\ =\ Perioden)\ \%\textgreater{}\%\ \ \ rename(Inflatie\ =\ Jaarmutatie.CPI....)\ \%\textgreater{}\%\ \ \ select(Jaartal,Inflatie)\ write.csv(inflatie,\ "data\ 2.0/Inflatie\_nieuw.csv",\ row.names\ =\ FALSE)}

\texttt{\{r.content-hidden\ when-format="pdf"\}\ \#Here\ we\ modify\ the\ dataset\ Welvaart\_van\_huishoudens\_\_kerncijfers\ welvaart\ \textless{}-\ welvaart\ \%\textgreater{}\%\ \ \ rename(leeftijd\_kostwinner\ =\ Kenmerken.van.huishoudens)\ \%\textgreater{}\%\ \ \ rename(Jaartal\ =\ Perioden)\ \%\textgreater{}\%\ \ \ rename(Gemiddeld\_gestandaardiseerd\_inkomen\_x1000\_euro\ =\ Gemiddeld.gestandaardiseerd.inkomen..1.000.euro.)\ head(welvaart)\ write.csv(welvaart,\ "data\ 2.0/Welvaart\_van\_huishoudens\_kerncijfers\_nieuw.csv",\ row.names\ =\ FALSE)}

\texttt{\{r.content-hidden\ when-format="pdf"\}\ \#Here\ we\ modify\ the\ dataset\ Bestaand\_koopwoningen\ verkoopprijzen\ Bestaande\_koopwoningen\_verkoopprijzen\_woningtype\ \textless{}-\ Bestaande\_koopwoningen\_verkoopprijzen\_woningtype\ \%\textgreater{}\%\ \ \ rename\ (Woningtype\ =\ Type.woning)\ \%\textgreater{}\%\ \ \ rename(\ Jaartal\ =\ Perioden)\ \%\textgreater{}\%\ \ \ rename(Aantal\_verkochte\_woningen\ =\ Verkochte.bestaande.koopwoningen.Verkochte.woningen..aantal.)\%\textgreater{}\%\ \ \ rename(\ Gemiddelde\_verkoopprijs\_euro\ =\ Gemiddelde.verkoopprijs..euro.)\%\textgreater{}\%\ \ \ select(Woningtype,\ Jaartal,\ Aantal\_verkochte\_woningen,\ Gemiddelde\_verkoopprijs\_euro)\ \ \ write.csv(Bestaande\_koopwoningen\_verkoopprijzen\_woningtype,\ "data\ 2.0/Bestaande\_koopwoningen\_verkoopprijzen\_woningtype\_nieuw.csv",\ row.names\ =\ FALSE)}

```\{r.content-hidden when-format=``pdf''\}

\#Start joining dataset inflatie\_nieuw and
Welvaart\_van\_huishoudens\_kerncijfers\_nieuw welvaart \textless-
read\_csv(``data
2.0/Welvaart\_van\_huishoudens\_kerncijfers\_nieuw.csv'') inflatie
\textless- read\_csv(``data 2.0/Inflatie\_nieuw.csv'') result \textless-
left\_join(welvaart, inflatie, by = ``Jaartal'')
write.csv2(result,``data 2.0/Samenvoeging1.csv'' , row.names = FALSE)

\begin{verbatim}


#Chapter 4: Housing availability

##4.1 Population growth and new houses built
```{r.content-hidden when-format="pdf"}
#Here we modify the dataset bevolkingsgroei to only keep the relevant information and use the right years.
bevolkingsgroei_nieuw <- bevolkingsgroei %>%
  rename( jong_volwassenen_20_tot_45_jaar = Bevolking.op.1.januari.Naar.leeftijd.20.tot.45.jaar..x.1.000.) %>%
  transmute(jong_volwassenen_20_tot_45_jaar) %>%
  mutate(Jaar = 2005:2019) %>%
  relocate(Jaar) %>%
  rename(Jaartal = Jaar)

write.csv(bevolkingsgroei_nieuw, "data 2.0/Bevolkingsgroei_nieuw.csv", row.names = FALSE)
\end{verbatim}

```\{r.content-hidden when-format=``pdf''\} \#Here we modify the dataset
Huizenprijzen\_Huizenvoorraad huizenvoorraad \textless- huizenvoorraad
\%\textgreater\% rename(Jaartal = Perioden) \%\textgreater\%
rename(Provincies = Regio.s) \%\textgreater\%
rename(Woningvoorraad\_1\_januari =
Bouwen.en.wonen.Woningvoorraad.Voorraad.op.1.januari..aantal.)\%\textgreater\%
rename(Nieuwbouw = Bouwen.en.wonen.Woningvoorraad.Nieuwbouw..aantal.)
\%\textgreater\% select(Jaartal, Provincies, Woningvoorraad\_1\_januari,
Nieuwbouw)

tabel\_per\_jaar \textless- huizenvoorraad \%\textgreater\%
group\_by(Jaartal) \%\textgreater\% summarise(
Woningvoorraad\_1\_januari = sum(Woningvoorraad\_1\_januari, na.rm =
TRUE), huizenvoorraad = sum(Nieuwbouw, na.rm = TRUE) )
write.csv(tabel\_per\_jaar, ``data 2.0/Huizenvoorraad\_nieuw.csv'',
row.names = FALSE)

\begin{verbatim}

```{r.content-hidden when-format="pdf"}
#Here we merge the datasets
bevolkingsgroei_nieuwbouw <- left_join(bevolkingsgroei_nieuw, tabel_per_jaar, by = "Jaartal")
write.csv(bevolkingsgroei_nieuwbouw, "data 2.0/bevolkingsgroei_nieuwbouw.csv")
\end{verbatim}

```\{r.content-hidden when-format=``pdf''\} \#Here we create a graph to
compare the amount of new houses built and the population growth of
young adults.

bevolkingsgroei\_nieuwbouw \textless- read.csv(``data
2.0/bevolkingsgroei\_nieuwbouw.csv'') \# d scale\_factor \textless-
max(bevolkingsgroei\_nieuwbouw\(jong_volwassenen_20_tot_45_jaar, na.rm = TRUE) /
                max(bevolkingsgroei_nieuwbouw\)huizenvoorraad, na.rm =
TRUE)

ggplot(bevolkingsgroei\_nieuwbouw, aes(x = Jaartal)) +

\# Left axis(population growth) geom\_line( aes(y =
jong\_volwassenen\_20\_tot\_45\_jaar, color = ``Young adults (20-45
years old)''), linewidth = 1.3 ) + geom\_point( aes(y =
jong\_volwassenen\_20\_tot\_45\_jaar, color = ``Young adults (20-45
years old)''), size = 2 ) + \# right axis(newly built homes) geom\_line(
aes(y = huizenvoorraad * scale\_factor, color = ``New houses built per
year''), linewidth = 1.3, linetype = ``dashed'' ) + geom\_point( aes(y =
huizenvoorraad * scale\_factor, color = ``New houses built per year''),
size = 2 ) + scale\_y\_continuous( \#left axis name = ``Young adults
(20-45 years old) x1000'', limits = c(3000, 6000), labels =
label\_number(big.mark = ``.'', decimal.mark = ``,''), breaks = c(3000,
4000, 5000, 6000),

\begin{verbatim}
#right axis
sec.axis = sec_axis(
  ~ . / scale_factor,
  name = "New houses built per year",
  breaks = c(40000, 50000, 60000, 70000, 80000),
  labels = label_number(big.mark = ".", decimal.mark = ",")
)
\end{verbatim}

) +

scale\_x\_continuous( breaks =
seq(min(bevolkingsgroei\_nieuwbouw\(Jaartal),
                 max(bevolkingsgroei_nieuwbouw\)Jaartal), by = 1) ) +

scale\_color\_manual( values = c( ``Young adults (20-45 years old)'' =
``blue'', ``New houses built per year'' = ``red'' ) ) +

labs( title = ``Growth of young adult population and the building of new
houses'', subtitle = ``Yearly deviations'', x = ``Year'', color = NULL )
+

theme\_minimal(base\_size = 14) +

theme( plot.title = element\_text(face = ``bold'', size = 18),
plot.subtitle = element\_text(size = 12), axis.text.x =
element\_text(size = 7), legend.position = ``top'' )

\begin{verbatim}
One thing to keep in mind beforehand is that New houses built, is a yearly growth variable and the population is the total population of the young adults in that year.

Taking this into account, we can see that there is a subtle downward trend in the population of young adults in the Netherlands. This trend is similar to other countries in Europe and is mainly caused by falling birth rates(Krzysztoszek, 2025). The other key thing to take away is that after 2009 the building of new houses declined rapidly until 2014. The shocks are of this brief decline are still felt years later. This decline is caused by the uncertainty after the credit crisis of 2008(Liefting & Berkhout, 2025). After 2014 we see an upward trend, but not by enough to compensate for the temporal decline.

Unfortunately we could not find numbers of population growth of this age group after 2019, so our analysis is limited in this area. Our expectations are however that the population will remain approximately equal. Our expectations of new houses being built are quite different. This is mainly based on the "stikstofcrisis". Thus, we expect less houses to built from 2019-2026(Liefting & Berkhout, 2025).

## 4.2 Housing per province
### trendgrafieken per provincie Huizenprijzen/voorraad

```{r.content-hidden when-format="pdf"}
df <- read_csv("data 2.0/Huizenprijzen_Huizenvoorraad_nieuw.csv") |> 
  clean_names() |> 
  mutate(provincie = str_remove(provincies, " \\(PV\\)"))

df_voorraad <- df |> 
  group_by(provincie) |> 
  arrange(jaartal) |> 
  mutate(index = woningvoorraad_1_januari / first(woningvoorraad_1_januari) * 100) |> 
  ungroup()

ggplot(df_voorraad, aes(jaartal, index, color = provincie)) +
  geom_line(linewidth = 0.9) +
  scale_x_continuous(breaks = seq(2005, 2025, 5)) +
  labs(title = "Growth in housing stock by province (2005 = 100)",
       x = "Year", y = "Index (2005 = 100)", color = "Province") +
  theme_minimal()

df_nat <- df |> 
  filter(jaartal < 2025) |> 
  group_by(jaartal) |> 
  summarise(nieuwbouw = sum(nieuwbouw, na.rm = TRUE))

ggplot(df_nat, aes(jaartal, nieuwbouw)) +
  geom_hline(yintercept = 100000, linetype = "dashed") +   # overheidsdoel
  geom_line(linewidth = 1, color = "steelblue") +
  geom_point() +
  scale_x_continuous(breaks = seq(2005, 2024, 3)) +
  scale_y_continuous(limits = c(0, 110000), labels = comma) +
  labs(title = "New housing construction remains below target",
       x = "Year", y = "Number of new-built homes") +
  theme_minimal()

df_intensiteit <- df |> 
  filter(jaartal < 2025) |> 
  mutate(per_1000 = nieuwbouw / woningvoorraad_1_januari * 1000)

ggplot(df_intensiteit, aes(jaartal, per_1000)) +
  geom_line(color = "darkorange") +
  facet_wrap(~ provincie) +
  scale_x_continuous(breaks = seq(2005, 2024, 5)) +
  labs(title = "Construction intensity by province",
       x = "", y = "New-built homes per 1,000") +
  theme_minimal()
\end{verbatim}

\#\#\#Explanation I put the housing data into three charts, because
housing stock and new construction are actually two different things.
The stock is the total number of homes that exist, and new construction
is what gets added each year. Putting those two in one chart with the
same index would be a bit misleading, since the stock rises slowly while
new construction jumps up and down a lot from year to year. In the first
chart I set the stock per province to 2005 = 100, which makes the
provinces comparable even though some are much bigger than others. In
the second chart I added up new construction nationally per year and put
a line at 100,000, the number the government wants to build each year.
One caveat: my figures only cover actual new construction, while that
target also counts things like offices turned into homes, so the line is
an indicative lower bound and not an exact match. The third chart shows
building intensity (new construction per 1,000 existing homes), which
corrects for how big a province is. I left 2025 out for new
construction, because that value is missing from the dataset, so I don't
draw conclusions about that year.

\#\#\#Conclusion Over time the housing stock grew steadily in every
province, but new construction tells the more interesting story. After a
peak around 2009 it dropped to a low of about 45,000 homes in 2014, and
only climbed back to around 74,000 by 2022, never reaching the target of
100,000. The timing of that dip lines up with the 2008 credit crisis,
but this is an association and not proof: the decline only set in with a
delay (building has a long pipeline), and the later recovery was also
shaped by other things like low interest rates and rising demand after
2020, so the crisis alone can't explain the whole pattern. It's also
important to say that my data only measures supply, not the shortage
itself. The link to demand comes from outside sources: the number of
households has been growing faster than the housing stock
(Rijksoverheid, n.d.), which left a statistical shortage of roughly
396,000 homes in 2025 (Volkshuisvesting Nederland, 2025). I can't show
that shortage from my own data alone, but putting the two together
suggests that for young adults the problem isn't only about prices, but
also about availability: too few homes were added to keep up with the
growing demand.

```\{r.content-hidden when-format=``pdf''\} library(jsonlite)

groei \textless- df \textbar\textgreater{} group\_by(provincie)
\textbar\textgreater{} arrange(jaartal) \textbar\textgreater{}
summarise(groei\_pct = (last(woningvoorraad\_1\_januari) /
first(woningvoorraad\_1\_januari) - 1) * 100)

geo \textless-
fromJSON(``\url{https://cartomap.github.io/nl/wgs84/provincie_2023.geojson}'',
simplifyVector = FALSE)

rows \textless- list(); g \textless- 0 for (f in geo\(features) {
  naam   <- f\)properties\(statnaam
  coords <- f\)geometry\(coordinates
  polys  <- if (f\)geometry\$type == ``Polygon'') list(coords) else
coords \# alles -\textgreater{} lijst van polygonen for (poly in polys)
for (ring in poly) \{ g \textless- g + 1 m \textless-
matrix(unlist(ring), ncol = 2, byrow = TRUE) \# ring -\textgreater{}
matrix van {[}lon, lat{]} rows{[}{[}length(rows) + 1{]}{]} \textless-
data.frame(provincie = naam, long = m{[},1{]}, lat = m{[},2{]}, group =
g) \} \} prov\_df \textless- do.call(rbind, rows)

\section{grenzen koppelen aan de
groeicijfers}\label{grenzen-koppelen-aan-de-groeicijfers}

plotdata \textless- prov\_df \textbar\textgreater{} left\_join(groei, by
= ``provincie'')

labels \textless- plotdata \textbar\textgreater{} group\_by(provincie)
\textbar\textgreater{} summarise(long = mean(long), lat = mean(lat),
.groups = ``drop'')

\section{choropleth tekenen}\label{choropleth-tekenen}

ggplot(plotdata, aes(long, lat, group = group, fill = groei\_pct)) +
geom\_polygon(color = ``white'', linewidth = 0.2) + geom\_text(data =
labels, aes(long, lat, label = provincie), inherit.aes = FALSE, size =
2.3, color = ``grey20'') + coord\_quickmap() + \# juiste verhouding
zonder sf scale\_fill\_viridis\_c(option = ``C'') + labs(title =
``Housing stock growth per province (2005--2025)'', caption = ``Source:
own analysis; province boundaries via CBS/cartomap'', fill = ``Growth
(\%)'') + theme\_void()

\begin{verbatim}

###Explanation
To show where in the country the growth happened, I mapped the total growth of the housing stock per province between 2005 and 2025. For each province I worked out the percentage change in the stock and used that to colour the map. The province borders come from a public CBS/cartomap file that I read in as a GeoJSON, so I didn't need any extra mapping software. One small caveat: the province labels sit at the average of each province's border points, which is only a rough centre, so for oddly shaped provinces the label can land slightly off. The colours and the actual figures aren't affected by that.

###Conclusion
The map shows a clear centre–periphery pattern. Growth is highest in the central provinces Flevoland (+34%) and Utrecht (+30%) and lowest in the ones on the edges, with Limburg in the far southeast (+13%) and the coastal provinces Zeeland and Fryslân (around +15%) at the bottom. This matters for the research question, because the stock grew fastest exactly in the central regions where demand from students and young workers is highest, while the quieter provinces on the edges grew the least. So the new supply wasn't added where it would help the most, which helps explain why housing stayed hard to find for young adults in precisely the areas they most want to live.


#Chapter 5 Housing affordability

```{r.content-hidden when-format="pdf"}
#Start joining dataset Samenvoeging1 and bruto_en_besteedbaar_inkomen_leeftijdsgroep_nieuw.csv
samenvoeging <- read_csv2("data 2.0/Samenvoeging1.csv")
samenvoeging2 <- left_join(brutoinkomen, samenvoeging, by = "Jaartal")
write.csv(samenvoeging2, "data 2.0/Samenvoeging2.csv" , row.names = FALSE)
head(samenvoeging2)
\end{verbatim}

\texttt{\{r.content-hidden\ when-format="pdf"\}\ \#Deleting\ a\ column\ that\ was\ mentioned\ twice\ samenvoeging2\ \textless{}-\ read.csv("data\ 2.0/Samenvoeging2.csv")\ samenvoeging2\ \textless{}-\ samenvoeging2\ \%\textgreater{}\%\ \ \ select\ (-leeftijd\_kostwinner)\ write.csv(samenvoeging2,\ "data\ 2.0/Final\_dataset.csv",\ row.names\ =\ FALSE)}
\#\# 5.1 Variables

The first, necessary variable added is the maximal nominal mortgage loan
that young adults can loan at the bank. The maximal nominal mortgage
loan is based on the LTI and the median gross income. The LTI
(Loan-to-Income) is the ratio between the gross income and the mortgage
loan. This ratio decides the maximum amount that can be borrowed, which
usually is around 4.5 times the gross income in the Netherlands
(Homefinance Hypotheken, n.d.). The choice to use the median gross
income over the average gross income is supported by the fact that the
mean usually is distorted by high extremes, whilst the median gives a
representative reflection of the ``typical'' young adult.

The Netherlands has one of the highest mortgage debts of Europe, which
states that the Dutch housing market is really dependent on credits
(DNB, n.d.). Because the Dutch housing market for buyers is
fundamentally credit-driven and reliant on the LTI, variables solely
addressing income are not enough to model the maximum purchasing
capacity of young adults. However, this nominal borrowing ceiling is the
variable that can demonstrate the maximum buying capacity of young
adults. Moreover, this variable can be rightfully compared to the
housing prices. All in all, the maximal nominal mortgage loan more or
less serves as the baseline for the study regarding housing
affordability.

\texttt{\{r.content-hidden\ when-format="pdf"\}\ af\_finaldataset\ \textless{}-\ read\_csv("data\ 2.0/Final\_dataset.csv")\ af\_newvariable\ \textless{}-\ af\_finaldataset\ \%\textgreater{}\%\ \ \ mutate(Max\_nominal\_mortgage\_loan\_1000.euro\ =\ Mediaan\_bruto\_inkomen\_1000.euro\ *\ 4.5)\ write.csv(af\_newvariable,\ "data\ 2.0/Final\_dataset\_2.csv",\ row.names\ =\ FALSE)}

Another variable that is valuable, especially in comparison with the
nominal maximum mortgage loan, is the real maximum mortgage loan. In
order to calculate this, we set the year 2024 as the base year and set
the adjustment factor to 1 in 2024. In this year, the nominal value is
equal to the real value. The adjustment factor quantifies how much
prices have risen with base year of 2024. The real maximum mortgage loan
is adjusted for inflation and clearly proves the money illusion and the
sudden inflation shocks. In order to create the extra column, it is
necessary to use the cumprod() function, as prices rise exponential.

\texttt{\{r.content-hidden\ when-format="pdf"\}\ af\_extranewvariable\ \textless{}-\ af\_newvariable\ \%\textgreater{}\%\ \ \ arrange\ (Jaartal)\ \%\textgreater{}\%\ \ \ mutate\ (price\_index\_2011\ =\ cumprod(1\ +\ (Inflatie/100))\ ,\ adjustment\_factor\_2024\ =\ price\_index\_2011\ {[}Jaartal\ ==\ 2024{]}\ /\ price\_index\_2011,\ Max\_real\_mortgage\_loan\_1000.euro\ =\ Max\_nominal\_mortgage\_loan\_1000.euro\ *\ adjustment\_factor\_2024\ )\ write.csv(af\_extranewvariable,\ "data\ 2.0/Final\_dataset\_with\_new\_variables.csv",\ row.names\ =\ FALSE)}

\section{Chapter 6 - Discussion}\label{chapter-6---discussion}

\subsection{6.1 Discuss your findings}\label{discuss-your-findings}

\section{Chapter 7 - Reproducibility}\label{chapter-7---reproducibility}

\subsection{7.1 Github repository link}\label{github-repository-link}

This is the link to our public repository:
\url{https://github.com/TimKnegt/Tutorial-group-1-1}

\subsection{7.2 Reference list}\label{reference-list}

\url{https://www.dnb.nl/actuele-economische-vraagstukken/woningmarkt/onze-hoge-hypotheekschulden-risico-s-en-oplossingen/}\strut \\
\url{https://www.homefinance.nl/hypotheek/informatie/lti-hypotheek/}

-CBS Statline. (2023, December 15). Centraal Bureau Voor De Statistiek.
Bevolking, huishoudens en bevolkingsontwikkeling; 1899-2019. Retrieved 5
june, 2026, from
\url{https://www.cbs.nl/nl-nl/cijfers/detail/37556?dl=34F0C}

-CBS Statline. (1995, season-01). Centraal Bureau voor de Statistiek.
Bestaande Koopwoningen; Verkoopprijzen; Woningtype; Prijsindex 2020=100.
Retrieved June 4, 2026, from
\url{https://opendata.cbs.nl/\#/CBS/nl/dataset/85791NED/table}

-CBS Statline. (2026, June 9). Centraal Bureau voor de Statistiek.
Jaarmutatie Consumentenprijsindex; Vanaf 1963. Retrieved June 10, 2026,
from \url{https://opendata.cbs.nl/\#/CBS/nl/dataset/70936ned/table}

-Krzysztoszek, A. (2025, June 12). Europe faces crisis as birth rates
drop, population ages. Pulse Z.
\url{https://www.pulse-z.eu/nl/de-europese-economie-loopt-gevaar-nu-het-geboortecijfer-daalt-en-de-bevolking-vergrijst/}

-Liefting, R., \& Berkhout, K. B. (2025, May 9). Zeven redenen waarom er
in Nederland te weinig huizen worden gebouwd -- Dutilhhuis. NRC.
\url{https://www.nrc.nl/nieuws/2025/05/09/zeven-redenen-waarom-er-in-nederland-te-weinig-huizen-worden-gebouwd-a4892667}

\end{document}
